\item \model{GLM-5}

\paragraph*{مبانی بهینه‌سازی و توجه تنک (\en{Optimization Foundations})}
مدل \model{GLM-5} با مقیاس ۷۴۴ میلیارد پارامتر کل (۴۰ میلیارد پارامتر فعال در هر توکن) و ۸۰ لایه ترنسفورمر، از سازوکار توجه تنک عمیق (\en{DeepSeek Sparse Attention - DSA}) برای مهار پیچیدگی محاسباتی در بافت‌های فراتر از ۱۲۸ هزار توکن استفاده می‌کند \cite{glm2026v5}. برای بهینه‌سازی وزن‌ها با بهینه‌ساز \en{Muon}، روش تفکیک هدهای توجه (\en{Muon Split}) اعمال شده است؛ در این روش، ماتریس‌های فرافکنی بزرگ $W^{UQ}, W^{UK}, W^{UV}$ به ماتریس‌های کوچک متناظر با هر هد تجزیه شده و تکرارهای متعامدسازی نیوتن-شولتز به صورت مستقل بر روی آن‌ها اجرا می‌شود تا نوسانات مقیاس گرادیان در لایه‌های پنهان کنترل گردد.

\paragraph*{تابع زیان استدلال با مهار انحراف توزیع (\en{Reasoning RL with IcePop})}
در مرحله آموزش استدلال، الگوریتم بهینه‌سازی سیاست نسبی گروهی (\en{GRPO}) \cite{shao2024deepseekmath} با سازوکار \en{IcePop} ترکیب شده تا ناهمخوانی میان توزیع استنتاج حین نمونه‌برداری مسیرها ($\pi_{\theta_{\text{old}}}^{\text{infer}}$) و توزیع آموزش حین اعمال گام‌های گرادیان ($\pi_{\theta_{\text{old}}}^{\text{train}}$) جبران گردد. 

نسبت ناهمخوانی آموزش-استنتاج در سطح توکن به‌صورت زیر تعریف می‌شود:
\begin{equation}
    \rho_{i,t} = \frac{\pi_{\theta_{\text{old}}}^{\text{train}}(y_{i,t} \mid x, y_{i,<t})}{\pi_{\theta_{\text{old}}}^{\text{infer}}(y_{i,t} \mid x, y_{i,<t})}
\end{equation}
عملگر فیلتر ناپیوستگی $\text{pop}(\rho_{i,t}, 1/\beta, \beta)$ مقادیر انحراف شدید را ماسک می‌کند:
\begin{equation}
    \text{pop}(\rho_{i,t}, 1/\beta, \beta) = \begin{cases} 
        \rho_{i,t}, & \text{if } \frac{1}{\beta} \le \rho_{i,t} \le \beta \\ 
        0, & \text{otherwise} 
    \end{cases}
\end{equation}
نسبت اهمیت گام بهینه‌سازی سیاست به فرم مرسوم زیر است:
\begin{equation}
    r_{i,t}(\theta) = \frac{\pi_\theta^{\text{train}}(y_{i,t} \mid x, y_{i,<t})}{\pi_{\theta_{\text{old}}}^{\text{train}}(y_{i,t} \mid x, y_{i,<t})}
\end{equation}
مزیت نرمال‌شده در گروه نمونه‌های به اندازه $G$:
\begin{equation}
    \hat{A}_{i,t} = \frac{R_i - \text{mean}(R_1, \dots, R_G)}{\text{std}(R_1, \dots, R_G)}
\end{equation}
با حذف جریمه واگرایی \en{KL} جهت تسریع در کشف فضاهای استدلال جدید، تابع زیان استدلال \model{GLM-5} عبارت است از:
\begin{multline}
    \mathcal{L}_{\text{Reasoning}}(\theta) = -\mathbb{E}_{x \sim \mathcal{D}, \, \{y_i\}_{i=1}^G \sim \pi_{\theta_{\text{old}}}^{\text{infer}}} \Bigg[ \frac{1}{G} \sum_{i=1}^G \frac{1}{|y_i|} \sum_{t=1}^{|y_i|} \text{pop}\left(\rho_{i,t}, \frac{1}{\beta}, \beta\right) \\
    \times \min \left( r_{i,t}\hat{A}_{i,t}, \; \text{clip}(r_{i,t}, 1-\epsilon_{\text{low}}, 1+\epsilon_{\text{high}})\hat{A}_{i,t} \right) \Bigg]
\end{multline}
که در آن $\beta = 2$، $\epsilon_{\text{low}} = 0.2$ و $\epsilon_{\text{high}} = 0.28$ تنظیم شده‌اند.

\paragraph*{یادگیری تقویتی ناهمگام عاملی (\en{Asynchronous Agentic RL})}
برای پشتیبانی از افق‌های زمانی بلند در حل مسائل مهندسی نرم‌افزار، چارچوب یادگیری تقویتی کاملاً ناهمگام طراحی شده است. از آنجا که مدل تولیدکننده مسیر با آهنگ متفاوتی نسبت به مدل یادگیرنده پیش می‌رود، ردیابی توزیع‌های گذشته ($\pi_{\theta_{\text{old}}}$) از طریق جایگزینی مستقیم نسبت احتمال ذخیره‌شده حین تولید رول‌اوت ($\pi_{\text{rollout}}$) انجام می‌شود:
\begin{equation}
    r_t(\theta) = \exp\left( \log \pi_\theta(a_t \mid s_t) - \log \pi_{\text{rollout}}(a_t \mid s_t) \right)
\end{equation}
ناحیه اعتماد از طریق تابع کالیبراسیون و برش توکنی دوطرفه زیر محدود می‌گردد:
\begin{equation}
    f(x; \epsilon_\ell, \epsilon_h) = \begin{cases} 
        x, & \text{if } 1 - \epsilon_\ell < x < 1 + \epsilon_h \\ 
        0, & \text{otherwise} 
    \end{cases}
\end{equation}
تابع هدف نهایی یادگیری عاملی ناهمگام عبارت است از:
\begin{equation}
    \mathcal{L}_{\text{Agent}}(\theta) = \mathbb{E}_t \left[ f(r_t(\theta), \epsilon_\ell, \epsilon_h) \hat{A}_t \log \pi_\theta(a_t \mid s_t) \right]
\end{equation}
مسیرهایی که فاصله نسخه‌ای آن‌ها از آستانه مجاز فراتر رود ($w' - w_0 > \tau$) یا ناشی از فروریزش محیط پایانه باشند، از محاسبه گرادیان حذف می‌گردند.

\paragraph*{تقطیر درون‌خطی میان‌مرحله‌ای (\en{On-Policy Cross-Stage Distillation})}
جهت ممانعت از فراموشی فاجعه‌بار پس از مراحل متعدد یادگیری تقویتی، در گام پایانی تقطیر متقابل با جایگزینی مزیت زیر در تابع زیان اجرا می‌شود:
\begin{equation}
    \hat{A}_{i,t} = \text{sg}\left[ \log \frac{\pi_{\theta_{\text{teacher}}}^{\text{infer}}(y_{i,t} \mid x, y_{i,<t})}{\pi_\theta^{\text{train}}(y_{i,t} \mid x, y_{i,<t})} \right]
\end{equation}

\paragraph*{دیاگرام ساختار یادگیری تقویتی و تطبیق توزیع در \model{GLM-5}}
\begin{center}
\begin{tikzpicture}[
    node distance=1.3cm and 1.5cm,
    font=\small,
    box/.style={rectangle, draw=blue!70!black, fill=blue!5, thick, rounded corners, minimum height=0.9cm, text centered, inner sep=5pt},
    op/.style={rectangle, draw=orange!80!black, fill=orange!10, thick, rounded corners, minimum height=0.9cm, text centered, inner sep=5pt},
    loss/.style={rectangle, draw=red!70!black, fill=red!8, thick, rounded corners, minimum height=0.9cm, text centered, inner sep=5pt},
    arrow/.style={-{Stealth[length=2.5mm]}, thick, draw=gray!80!black}
]
    \node[box] (rollout) {نمونه‌برداری ناهمگام $\pi_{\text{rollout}}$};
    \node[op, right=of rollout] (ratio) {محاسبه نسبت اهمیت توکنی $r_t(\theta)$};
    \node[op, right=of ratio] (pop) {اعمال برش دوطرفه $f(r_t)$ و فیلتر \en{IcePop}};
    \node[loss, below=of ratio] (loss) {زیان بهینه‌سازی گرادیان سیاست عاملی $\mathcal{L}_{\text{Agent}}$};
    \node[box, left=of loss] (distill) {تقطیر میان‌مرحله‌ای با گرادیان شکاف لاجیت};
    
    \draw[arrow] (rollout) -- (ratio);
    \draw[arrow] (ratio) -- (pop);
    \draw[arrow] (pop) |- (loss);
    \draw[arrow] (distill) -- (loss);
\end{tikzpicture}
\end{center}